\usepackage{amssymb}
\usepackage{paralist}
\usepackage{bbm}
\usepackage{mathtools}
\usepackage{tikz}
\usepackage{natbib}
\usepackage{cleveref}

\newcommand{\eol}{\$}
\renewcommand{\O}{\mathrm{O}}
\newcommand{\tildeO}{\widetilde{\mathrm{O}}}
\DeclareMathOperator{\sgn}{sgn}

% Floor and ceiling functions.
% \DeclarePairedDelimiter\ceil{\lceil}{\rceil}
% \DeclarePairedDelimiter\floor{\lfloor}{\rfloor}

% Norms and stuff.
\DeclarePairedDelimiter\abs{\lvert}{\rvert}%
\DeclarePairedDelimiter\norm{\lVert}{\rVert}%

% Argmin and argmax
% \DeclareMathOperator*{\argmin}{argmin}
% \DeclareMathOperator*{\argmax}{argmax}

% Swap the definition of \abs* and \norm*, so that \abs
% and \norm resizes the size of the brackets, and the 
% starred version does not.
\makeatletter
\let\oldabs\abs
\def\abs{\@ifstar{\oldabs}{\oldabs*}}
%
\let\oldnorm\norm
\def\norm{\@ifstar{\oldnorm}{\oldnorm*}}
\makeatother

\newcommand{\draftcomment}[3]{{\textcolor{#3}{[#2: \emph{#1}]}}}
% % \renewcommand{\draftcomment}[3]{}  % uncomment for submission
\newcommand{\will}[1]{\draftcomment{#1}{\textsc{wm}}{orange}}
\newcommand{\ashish}[1]{\draftcomment{#1}{\textsc{as}}{brown}}

\let\polishL\L  % preserve the original meaning of \L before redefining it below

\renewcommand\P{\mathsf P}
\renewcommand\L{\mathsf L}
\newcommand\NL{\mathsf{NL}}
\newcommand{\TC}{\mathsf{TC}}
\newcommand{\NC}{\mathsf{NC}}
\newcommand\PSPACE{\mathsf{PSPACE}}
\newcommand\FOM{\mathsf{FO(M)}}

\newcommand{\CoT}{\mathsf{CoT}}
\newcommand{\calT}{\mathcal T}
\newcommand\TIME{\mathsf{TIME}}
\newcommand\tildeTIME{\widetilde{\mathsf{TIME}}}
\newcommand{\SPACE}{\mathsf{SPACE}}
\newcommand{\CSL}{\mathsf{CSL}}

\newcommand{\poly}{\mathrm{poly}}

% Layer norm notation.
\newcommand{\lnorm}{\mathsf{layer\_norm}}
\newcommand{\mlnorm}{\mathsf{multi\_layer\_norm}}
\newcommand{\slnorm}{\mathsf{proj\_layer\_norm}}
\newcommand{\hash}{\phi}

\usepackage{amsthm}
\usepackage{thmtools}
\usepackage{thm-restate}
\newtheorem{theorem}{Theorem}
\newtheorem{proposition}{Proposition}
\newtheorem{lemma}{Lemma}
\newtheorem{corollary}{Corollary}[theorem]
\newtheorem{lemcorollary}{Corollary}[lemma]
\theoremstyle{definition}
\newtheorem{definition}{Definition}
\newtheorem{example}{Example}